\begin{remark}[Reader precedence for historical release snapshots inside the active line]
The next two synchronized-status blocks reconstruct older public-stage status only for navigation and burden genealogy. They do not overrule the current burden ledger, the current theorematic status, or the current tool/QAM/decoder routing of the active line. Whenever one of those historical snapshots sounds flatter or stronger than the normalized present reading, the normalized present reading governs.
\end{remark}

\begin{remark}[v3.23 synchronized status reading for the optical approximation layer]
\label{rem:v323-synchronized-status-reading-for-the-optical-approximation-layer}
The historical public v3.23.0 release line no longer permits a flat reading in which the new optical sector is either ``already a full optical theory'' or ``still only a loose analogy.'' After the first reduced prism-readout lemma, the residual-diagnostic layer, the carrier-side fit/calibration/validation chain, the reconstruction/correction steps, and the bundled optical approximation package, the optical layer now has its own staged status:
\begin{enumerate}[leftmargin=2em,label=(\roman*)]
  \item a \emph{strong global optical claim} remains open;
  \item a \emph{reduced retained-carrier optical sector} is already available;
  \item a \emph{first-order optical workflow} is already available;
  \item a \emph{bundled first-order optical package} is already available as an active citation surface.
\end{enumerate}
Thus the open-front layer must now distinguish between the still-open stronger global optical theory and the already stabilized first-order optical approximation layer.
\end{remark}

\begin{proposition}[v3.23 optical/open-question status synchronization principle]
\label{prop:v323-optical-open-question-status-synchronization-principle}
At the present theorematic and retained-carrier level, the following synchronized status reading holds for the new optical line:
\begin{enumerate}[leftmargin=2em,label=(\roman*)]
  \item the still-open strong claim is a full global optical theory of the crystal--shadow architecture, including non-reduced anisotropic branches, stronger inversion, and a later convergence theory for higher optical refinements;
  \item what is already available is a reduced prism-readout sector on one retained carrier surface with ordinary wavelength ordering;
  \item what is already available as first-order workflow is the chain consisting of residual diagnostics, carrier-side fit ansatz, spectral calibration, overdetermined validation, effective-index reconstruction, carrier-profile correction, and reduced stability/stopping control;
  \item what is already available as a single active surface is the bundled optical approximation package of Proposition~\ref{prop:first-bundled-optical-approximation-package};
  \item therefore the optical sector should no longer be counted among the fully open foundational fronts of the document, but as a partially stabilized first-order approximation layer with a still-open stronger global continuation.
\end{enumerate}
\end{proposition}

\begin{proof}
Clause (i) records exactly the stronger optical ambitions that the current line has not yet proved: unrestricted anisotropic branching, stronger inversion, and a genuine global convergence theory are still absent. Clause (ii) is supplied by the reduced prism-readout lemma. Clause (iii) is the cumulative content of the residual spectral-diagnostic lemma, the carrier-side fit principle, the calibration principle, the overdetermined consistency principle, the effective-index reconstruction principle, the carrier-profile correction principle, and the reduced stability principle. Clause (iv) is exactly the content of the bundled optical approximation package. Hence the optical line already has a stabilized first-order approximation layer even though the stronger global optical theory remains open.
\end{proof}

\begin{corollary}[Use rule for the synchronized v3.23 optical status reading]
\label{cor:use-rule-for-the-synchronized-v323-optical-status-reading}
Whenever a later argument in the present public v3.23.0 release line needs to classify whether the optical sector is already available for active use, it is sufficient to cite Proposition~\ref{prop:v323-optical-open-question-status-synchronization-principle}. A finer reopening is required only when the argument depends on a stronger global optical theory beyond the reduced retained-carrier approximation layer.
\end{corollary}

\begin{proof}
This is the operational reformulation of Proposition~\ref{prop:v323-optical-open-question-status-synchronization-principle}. The proposition already isolates the precise point at which the optical line ceases to count as merely a loose heuristic while still falling short of a full global optical theory.
\end{proof}


\begin{remark}[Historical v3.22/v3.23 synchronized status reading for the upgraded closure layer]
\label{rem:v322-synchronized-status-reading-upgraded-closure-layer}
The historical public v3.23.0 release line no longer permits a flat reading in which the upgraded closure/interface blocks are simply either ``open'' or ``proved.'' After the public v3.20 theorem-level shadow package, the public v3.21 interface-payoff layer, and the present v3.22 partial-closure steps, their status is now genuinely four-tiered:
\begin{enumerate}[leftmargin=2em,label=(\roman*)]
  \item a \emph{strong deferred claim} remains open;
  \item a \emph{theorem-level shadow} is already available;
  \item an \emph{interface-level operational use} is already available;
  \item for certain carrier-side conclusions, a \emph{restricted partial-closure / lift-back consequence} is already available.
\end{enumerate}
Thus the open-front layer must now distinguish not only between strong and weak forms, but also between theorematic shadow, interface use, and restricted closure consequence.
\end{remark}

\begin{proposition}[v3.22 hypothesis/open-question status synchronization principle]
\label{prop:v322-hypothesis-open-question-status-synchronization-principle}
At the present theorematic and interface level, the following synchronized status reading holds:
\begin{enumerate}[leftmargin=2em,label=(\roman*)]
  \item for Hypothesis~\ref{hyp:hc-meta-generator}, the still-open strong claim is the existence of one explicit compensator-centered generating equation $\mathfrak E_{\mathrm{HC}}=0$, while the theorem-level shadow of a common closure-stable carrier/projection architecture is already available, already usable on the terminal interface, and now participates in the restricted lift-back layer whenever only retained-carrier conclusions are drawn;
  \item for Hypothesis~\ref{hyp:frame-dependent-defect-readout}, the still-open strong claim is the existence of a deeper explicit compensator generator $\mathbf H_C$, while the theorem-level shadow separating visible defect vanishing from carrier extinction is already available, already usable on the terminal interface, and now participates in the no-second-carrier / restricted lift-back layer;
  \item for Hypothesis~\ref{hyp:fluid-readout}, the still-open strong claim is the existence of a concrete coarse-graining map to fluid-type evolution, while the theorem-level shadow demoting fluid-like behavior to a later coarse-grained readout is already available, already usable on the terminal interface, and now participates negatively in the restricted lift-back layer by forbidding replacement of the retained carrier;
  \item for Hypothesis~\ref{hyp:bidirectional-probe-lensing}, the still-open strong claim is a full universal geometric lensing law, while the theorem-level shadow of common-carrier/two-probe coupling is already available, already usable on the terminal interface, and now participates in the no-second-carrier / partial-closure layer;
  \item for the OP~3/OP~5 line, the still-open strong claim is a full global closure theorem, while the historical public v3.23.0 release line already contained more than a weak bridge: it had a first partial closure theorem together with a restricted lift-back to retained carrier-layer structure.
\end{enumerate}
Therefore these upgraded blocks should no longer be counted among the \emph{fully open foundational fronts} of the document. Only their strong deferred forms remain fully open; their weaker theorematic, interface-level, and partial-closure forms are already part of the active line.
\end{proposition}

\begin{proof}
Each clause follows by combining the public v3.20 theorem-level shadows, the public v3.21 interface-level application principles, and where applicable the new v3.22 partial-closure consequences. For the HC meta-field generator hypothesis, this combines the common-carrier/projection shadow and its interface-level use with the present restricted lift-back layer. For frame-dependent defect readout, this combines the defect-shadow theorem, the interface-level defect-face vanishing principle, and the present no-second-carrier/lift-back layer. For fluid-type readout, this combines the fluid-shadow theorem, the interface-level fluid demotion principle, and the present lift-back restriction that forbids carrier replacement. For bidirectional probe lensing, this combines the probe-lensing shadow theorem, the interface-level probe-coupling principle, and the present no-second-carrier/partial-closure layer. Finally, the OP~3/OP~5 line is upgraded from the v3.21 bridge test to the present v3.22 partial closure theorem and restricted lift-back theorem. Hence that historical stage already supported a synchronized four-tier status reading.
\end{proof}

\begin{corollary}[Use rule for the synchronized v3.22 status reading]
\label{cor:use-rule-for-the-synchronized-v322-status-reading}
Whenever a later argument refers back to the historical public v3.23.0 release line and needs to classify whether one of the upgraded closure/interface blocks was still truly open at that stage, it is sufficient to cite Proposition~\ref{prop:v322-hypothesis-open-question-status-synchronization-principle}. A finer reopening is required only when the argument depends on the still-open strong form rather than on the already stabilized theorematic, interface-level, or restricted partial-closure form.
\end{corollary}

\begin{proof}
This is the operational reformulation of Proposition~\ref{prop:v322-hypothesis-open-question-status-synchronization-principle}. The proposition already isolates the precise point at which the upgraded blocks cease to count as fully open foundational fronts.
\end{proof}


\begin{remark}[v3.21 synchronized status reading for the upgraded hypothesis layer]
\label{rem:v321-synchronized-status-reading-upgraded-hypothesis-layer}
The present v\docversion\ release line no longer permits a flat reading in which the upgraded hypothesis blocks are simply either ``fully open'' or ``already proved.'' After the released v3.20 theorem-level shadow package and the historical v3.21 interface payoff layer, their status is now genuinely three-tiered:
\begin{enumerate}[leftmargin=2em,label=(\roman*)]
  \item a \emph{strong deferred claim} remains open;
  \item a \emph{theorem-level shadow} is already available;
  \item an \emph{interface-level operational use} is already available on the terminal surface.
\end{enumerate}
Thus the open-front layer must now distinguish between still-open strong forms and already stabilized weaker forms.
\end{remark}

\begin{proposition}[v3.21 hypothesis/open-question status synchronization principle]
\label{prop:v321-hypothesis-open-question-status-synchronization-principle}
At the present theorematic and interface level, the following synchronized status reading holds:
\begin{enumerate}[leftmargin=2em,label=(\roman*)]
  \item for Hypothesis~\ref{hyp:hc-meta-generator}, the still-open strong claim is the existence of one explicit compensator-centered generating equation $\mathfrak E_{\mathrm{HC}}=0$, while the theorem-level shadow of a common closure-stable carrier/projection architecture is already available and already usable on the terminal interface;
  \item for Hypothesis~\ref{hyp:frame-dependent-defect-readout}, the still-open strong claim is the existence of a deeper explicit compensator generator $\mathbf H_C$, while the theorem-level shadow separating visible defect vanishing from carrier extinction is already available and already usable on the terminal interface;
  \item for Hypothesis~\ref{hyp:fluid-readout}, the still-open strong claim is the existence of a concrete coarse-graining map to fluid-type evolution, while the theorem-level shadow demoting fluid-like behavior to a later coarse-grained readout is already available and already usable on the terminal interface;
  \item for Hypothesis~\ref{hyp:bidirectional-probe-lensing}, the still-open strong claim is a full universal geometric lensing law, while the theorem-level shadow of common-carrier/two-probe coupling is already available and already usable on the terminal interface.
\end{enumerate}
Therefore these four blocks should no longer be counted among the \emph{fully open foundational fronts} of the document. Only their strong deferred forms remain open; their weaker theorematic and interface-operational forms are already part of the active line.
\end{proposition}

\begin{proof}
Each clause follows by combining one theorem-level shadow result with its corresponding interface-level application principle. For Hypothesis~\ref{hyp:hc-meta-generator}, this is the theorem-level shadow of the HC meta-field generator together with the interface-level common-carrier/projection principle. For Hypothesis~\ref{hyp:frame-dependent-defect-readout}, this is the theorem-level shadow of frame-dependent defect readout together with the interface-level defect-face vanishing principle. For Hypothesis~\ref{hyp:fluid-readout}, this is the theorem-level shadow of fluid-type readout together with the interface-level fluid-readout demotion principle. For Hypothesis~\ref{hyp:bidirectional-probe-lensing}, this is the theorem-level shadow of probe-lensing together with the interface-level probe-coupling principle. Hence the four upgraded blocks now share one synchronized three-tier status reading: strong deferred form, theorem-level shadow, and interface-level operational use.
\end{proof}

\begin{corollary}[Use rule for the synchronized v3.21 status reading]
\label{cor:use-rule-for-the-synchronized-v321-status-reading}
Whenever a later argument in the present v\docversion\ release line needs to classify whether one of the four upgraded hypothesis blocks is still truly open, it is sufficient to cite Proposition~\ref{prop:v321-hypothesis-open-question-status-synchronization-principle}. A finer reopening is required only when the argument depends on the still-open strong form rather than on the already stabilized theorematic or interface-level shadow.
\end{corollary}

\begin{proof}
This is the operational reformulation of Proposition~\ref{prop:v321-hypothesis-open-question-status-synchronization-principle}. The proposition already isolates the precise point at which the upgraded blocks cease to count as fully open foundational fronts.
\end{proof}


